\documentclass[11pt,a4paper]{article}
\usepackage[margin=28mm]{geometry}
\usepackage{amsmath,amssymb}
\usepackage{iftex}
\usepackage{kotex}
\ifPDFTeX
  \usepackage[T1]{fontenc}
\else
  \IfFontExistsTF{Malgun Gothic}{\setmainhangulfont{Malgun Gothic}}{}
\fi

\title{Notes on Continuous Optimization}
\author{Research notes}
\date{November 2026}

\begin{document}
\maketitle

\begin{abstract}
Continuous optimization offers a simple language for describing how systems
improve. These notes introduce the objective function, the geometry of a
gradient, and the small steps that connect a local decision to a better
solution.
\end{abstract}

이 문서는 UTF-8로 저장되며, XeLaTeX와 kotex로 한글을 조판합니다.

\section{Introduction}
Many questions in science and engineering can be expressed as a search for
the best possible choice. We choose a model, define what it means to
improve, and let a sequence of small updates guide us.

The central idea is simple: move in a direction that reduces the objective.
For a differentiable function, its gradient tells us which direction to take.

\section{Problem formulation}
Let $f : \mathbb{R}^n \to \mathbb{R}$ be a differentiable function.
Our task is to find a point that minimizes its value:
\begin{equation}
  x^\star = \operatorname*{arg\,min}_{x \in \mathbb{R}^n} f(x).
\end{equation}

\subsection{Gradient descent}
Starting from an initial estimate $x_0$, we repeat the update
\begin{equation}
  x_{k+1} = x_k - \eta\,\nabla f(x_k),
\end{equation}
where $\eta > 0$ is the learning rate. A small enough step makes each
local decision a useful part of the larger search.

\section{Conclusion}
Optimization connects a mathematical objective to a practical procedure.
The quality of a solution depends on the model, the step size, and the
geometry of the problem. These simple ideas form the foundation for
many modern learning algorithms.

\end{document}
